% Generated by generate.py: do not edit.
\chapter{Part I: Foundations}
\section{Tome 0: Back in Shape}
\subsection*{1.\ Sums, Binomials and Induction}
\begin{results}
\item[1.1] Proposition (induction principles)
\item[1.3] Proposition (telescoping sums)
\item[1.6] Proposition (Vandermonde)
\item[1.8] Theorem (Newton's forward difference formula)
\item[1.9] Corollary (sums of polynomials)
\item[1.10] Proposition (counting paths)
\end{results}
\subsection*{2.\ Complex Numbers and Roots of Unity}
\begin{results}
\item[2.1] Proposition (exponential form)
\item[2.3] Proposition (roots of unity)
\item[2.4] Proposition (orthogonality)
\item[2.5] Theorem (roots-of-unity filter)
\item[2.7] Proposition (equilateral triangles)
\end{results}
\subsection*{3.\ Polynomials}
\begin{results}
\item[3.1] Proposition (Euclidean division)
\item[3.2] Proposition (roots)
\item[3.3] Theorem (d'Alembert--Gauss)
\item[3.4] Proposition (Vi\`ete's relations)
\item[3.5] Theorem (Lagrange)
\item[3.7] Proposition (rational roots)
\item[3.8] Proposition (Gauss' lemma)
\item[3.9] Theorem (Eisenstein's criterion)
\item[3.10] Proposition (Newton's identities)
\item[3.11] Theorem (fundamental theorem of symmetric polynomials)
\item[3.12] Theorem (Descartes' rule of signs)
\item[3.13] Theorem (Gauss--Lucas)
\end{results}
\subsection*{4.\ Inequalities}
\begin{results}
\item[4.1] Proposition (convexity tests)
\item[4.2] Theorem (Jensen)
\item[4.3] Theorem (AM--GM)
\item[4.4] Theorem (Cauchy--Schwarz)
\item[4.5] Proposition (rearrangement)
\item[4.6] Theorem (power means)
\item[4.7] Proposition (Chebyshev's sum inequality)
\item[4.8] Theorem (Minkowski)
\end{results}
\subsection*{5.\ Sequences, Series and Limits}
\begin{results}
\item[5.1] Proposition (convergence criteria)
\item[5.2] Theorem (Ces\`aro)
\item[5.3] Proposition (series of non-negative terms)
\item[5.4] Proposition (signed series)
\item[5.5] Proposition (Euler's constant)
\item[5.6] Theorem (Stirling)
\item[5.7] Theorem (intermediate values)
\item[5.8] Theorem (extreme values)
\item[5.9] Theorem (Heine: uniform continuity)
\item[5.10] Theorem (Rolle, and the mean value theorem)
\item[5.11] Proposition (Abel's summation by parts)
\item[5.12] Theorem (Dirichlet's test)
\end{results}
\subsection*{6.\ Integrals and Taylor Formulas}
\begin{results}
\item[6.1] Proposition (integration toolkit)
\item[6.2] Theorem (Taylor formulas)
\item[6.4] Theorem (Stirling's constant)
\item[6.5] Theorem (l'H\^opital's rule)
\item[6.6] Proposition (the Gamma function)
\item[6.7] Proposition (midpoint and trapezoid rules)
\item[6.8] Proposition (Stirling, one term further)
\end{results}
\subsection*{7.\ Linear Algebra, Revisited}
\begin{results}
\item[7.1] Lemma (sign of a permutation)
\item[7.2] Theorem (determinant)
\item[7.3] Proposition (rank--nullity)
\item[7.4] Proposition (Vandermonde)
\item[7.5] Proposition (characteristic polynomial)
\item[7.6] Theorem (diagonalisability)
\item[7.7] Theorem (spectral theorem)
\item[7.9] Theorem (Gram--Schmidt and orthogonal projection)
\item[7.10] Theorem (Gram determinant)
\item[7.11] Proposition (adjugate, Cramer's rule)
\item[7.12] Theorem (Cayley--Hamilton)
\end{results}
\subsection*{8.\ Counting and First Probability}
\begin{results}
\item[8.1] Proposition (basic counts)
\item[8.2] Theorem (inclusion--exclusion)
\item[8.3] Proposition (derangements)
\item[8.4] Proposition (first rules)
\item[8.6] Proposition (products and sums)
\item[8.7] Proposition (the classical laws)
\item[8.8] Theorem (Markov and Chebyshev)
\item[8.9] Theorem (weak law of large numbers)
\end{results}
\section{Tome 1: The Integers}
\subsection*{1.\ Divisibility and the Euclidean Algorithm}
\begin{results}
\item[1.1] Theorem (Euclidean division)
\item[1.4] Theorem (B\'ezout)
\item[1.5] Theorem (Gauss' lemma)
\item[1.6] Theorem (fundamental theorem of arithmetic)
\item[1.8] Theorem (Lam\'e)
\end{results}
\subsection*{2.\ Primes}
\begin{results}
\item[2.1] Theorem (Euclid)
\item[2.2] Theorem (Euler)
\item[2.4] Theorem (Chebyshev-type bounds)
\item[2.5] Theorem (Bertrand's postulate; Chebyshev 1852, proof by Erd\H{o}s 1932)
\end{results}
\subsection*{3.\ Congruences}
\begin{results}
\item[3.3] Theorem (Euler, Fermat)
\item[3.4] Theorem (Wilson)
\item[3.5] Theorem (Chinese remainder theorem)
\item[3.6] Proposition (linear congruences)
\item[3.7] Theorem (Hensel's lemma)
\end{results}
\subsection*{4.\ Arithmetic Functions}
\begin{results}
\item[4.4] Theorem (M\"obius inversion)
\item[4.5] Proposition (Gauss)
\item[4.6] Theorem (Euclid--Euler)
\item[4.7] Theorem (Dirichlet)
\end{results}
\subsection*{5.\ Primitive Roots and the Multiplicative Group}
\begin{results}
\item[5.3] Theorem (Gauss)
\end{results}
\subsection*{6.\ Quadratic Residues and Reciprocity}
\begin{results}
\item[6.2] Theorem (Euler's criterion)
\item[6.3] Theorem (Gauss' lemma)
\item[6.5] Theorem (quadratic reciprocity; Gauss, 1796)
\item[6.6] Lemma (Eisenstein)
\end{results}
\subsection*{7.\ Sums of Squares}
\begin{results}
\item[7.1] Lemma (Thue)
\item[7.2] Theorem (Fermat, 1640; Euler, 1749)
\item[7.4] Theorem (Lagrange, 1770)
\item[7.5] Theorem (Gauss)
\end{results}
\subsection*{8.\ Continued Fractions and Pell's Equation}
\begin{results}
\item[8.3] Lemma (best approximation)
\item[8.4] Theorem (Legendre)
\item[8.5] Theorem (Euler, Lagrange)
\item[8.7] Theorem (Lagrange, 1768)
\end{results}
\subsection*{9.\ Diophantine Equations}
\begin{results}
\item[9.2] Theorem (Fermat)
\item[9.5] Theorem (Lebesgue, 1869)
\item[9.6] Theorem (Markov, 1879)
\end{results}
\subsection*{10.\ Gaussian and Eisenstein Integers}
\begin{results}
\item[10.2] Proposition (Euclidean division)
\item[10.5] Theorem (Jacobi)
\item[10.8] Theorem (Euler; Gauss)
\end{results}
\subsection*{11.\ Ideals: a First Look at Algebraic Number Theory}
\begin{results}
\item[11.4] Theorem (Dirichlet's unit theorem, real quadratic case)
\item[11.6] Proposition (Kummer's ideal factors of $6$)
\item[11.10] Theorem (Dedekind)
\end{results}
\subsection*{12.\ The \texorpdfstring{$p$}{p}-adic Numbers}
\begin{results}
\item[12.7] Theorem (Hensel's lemma, Newton's form)
\item[12.8] Corollary (squares in $\Q_p$)
\item[12.9] Corollary (roots of unity)
\item[12.10] Theorem (Ostrowski, 1916)
\end{results}
\section{Tome 2: Structures}
\subsection*{1.\ Groups as Symmetries}
\begin{results}
\item[1.2] Proposition (first rules)
\item[1.4] Proposition (the dihedral group)
\item[1.5] Proposition (rotations of the cube)
\item[1.8] Proposition (order of an element)
\item[1.10] Theorem (Lagrange)
\item[1.13] Theorem (cyclic groups)
\end{results}
\subsection*{2.\ Morphisms and Quotients}
\begin{results}
\item[2.6] Theorem (quotient group)
\item[2.7] Theorem (first isomorphism theorem)
\item[2.9] Theorem (second and third isomorphism theorems, correspondence)
\item[2.11] Theorem (Chinese remainder theorem, group version)
\item[2.12] Proposition (internal direct product)
\end{results}
\subsection*{3.\ Group Actions and Counting}
\begin{results}
\item[3.3] Theorem (Cayley)
\item[3.5] Theorem (orbit--stabiliser)
\item[3.6] Theorem (Burnside's lemma)
\item[3.7] Theorem (P\'olya's counting theorem)
\item[3.9] Theorem (class equation)
\item[3.11] Theorem (Cauchy)
\end{results}
\subsection*{4.\ The Symmetric Group}
\begin{results}
\item[4.2] Proposition (cycle decomposition)
\end{results}
\subsection*{5.\ Sylow Theorems and Finite Groups}
\begin{results}
\item[5.3] Theorem (Sylow I)
\item[5.4] Theorem (Sylow II and III)
\item[5.8] Theorem (groups of order $pq$)
\end{results}
\subsection*{6.\ Rings and Ideals}
\begin{results}
\item[6.5] Proposition (quotient rings)
\item[6.9] Theorem (Chinese remainder theorem for rings)
\item[6.15] Lemma (Gauss)
\item[6.17] Theorem (Gauss)
\end{results}
\subsection*{7.\ Modules over a Principal Ideal Domain}
\begin{results}
\item[7.4] Theorem (Smith normal form)
\item[7.6] Theorem (structure of finitely generated modules over a PID)
\item[7.7] Corollary (finitely generated abelian groups)
\item[7.10] Theorem (rational canonical form; Frobenius)
\item[7.11] Corollary (Cayley--Hamilton)
\item[7.12] Theorem (Jordan normal form)
\end{results}
\subsection*{8.\ Field Extensions}
\begin{results}
\item[8.2] Theorem (tower law)
\item[8.5] Proposition (minimal polynomial)
\item[8.6] Theorem (simple extensions)
\item[8.9] Theorem (Kronecker)
\item[8.14] Theorem (uniqueness of splitting fields)
\item[8.17] Theorem (Wantzel)
\item[8.18] Corollary (the classical problems)
\end{results}
\subsection*{9.\ Finite Fields}
\begin{results}
\item[9.2] Lemma (Frobenius)
\item[9.6] Theorem (Galois, Moore)
\item[9.10] Theorem (Gauss)
\item[9.12] Theorem (Berlekamp--Welch decoding)
\end{results}
\subsection*{10.\ Galois Theory}
\begin{results}
\item[10.5] Lemma (Artin)
\item[10.6] Theorem (fundamental theorem of Galois theory)
\item[10.10] Theorem (Gauss, Wantzel)
\item[10.13] Theorem (Galois)
\item[10.14] Theorem (Abel, Ruffini)
\item[10.15] Lemma (Dedekind)
\item[10.16] Proposition (Kummer)
\item[10.17] Theorem (Galois' criterion)
\end{results}
\subsection*{11.\ Representations and Characters}
\begin{results}
\item[11.3] Theorem (Maschke)
\item[11.4] Lemma (Schur)
\item[11.8] Theorem (orthogonality of characters)
\item[11.11] Corollary (column orthogonality)
\item[11.14] Theorem (Fourier analysis on a finite abelian group)
\end{results}
\subsection*{12.\ A Glimpse Beyond: Quaternions and Rotations}
\begin{results}
\item[12.3] Lemma (Euler)
\end{results}
\section{Tome 3: Chance}
\subsection*{1.\ Counting and Paradoxes}
\begin{results}
\item[1.11] Lemma (reflection principle)
\item[1.12] Theorem (Bertrand's ballot theorem, 1887)
\item[1.13] Theorem (Buffon)
\end{results}
\subsection*{2.\ Discrete Random Variables}
\begin{results}
\item[2.4] Proposition (transfer formula)
\item[2.6] Proposition (Cauchy--Schwarz)
\item[2.11] Theorem (Poisson's law of rare events)
\item[2.12] Theorem (Feller's coupling)
\end{results}
\subsection*{3.\ Generating Functions}
\begin{results}
\item[3.8] Theorem (extinction)
\end{results}
\subsection*{4.\ Inequalities, Large Numbers and Entropy}
\begin{results}
\item[4.2] Theorem (weak law of large numbers)
\item[4.4] Lemma (Heine)
\item[4.5] Theorem (Weierstrass, Bernstein's proof)
\item[4.6] Theorem (Chernoff bounds)
\item[4.7] Theorem (large deviations for coin flips)
\item[4.8] Lemma (Hoeffding)
\item[4.9] Theorem (Hoeffding's inequality)
\item[4.11] Theorem (Borel--Cantelli lemmas)
\item[4.12] Corollary (strong law for bounded variables)
\item[4.14] Theorem (Jensen)
\item[4.16] Proposition (Gibbs' inequality)
\item[4.18] Proposition (Kraft's inequality)
\item[4.19] Theorem (Shannon's source coding theorem)
\item[4.21] Theorem (asymptotic equipartition property)
\item[4.22] Corollary (fixed-length compression)
\end{results}
\subsection*{5.\ Continuous Probability Done Properly}
\begin{results}
\item[5.1] Theorem (Vitali, 1905)
\item[5.4] Theorem (Dynkin's $\pi$-$\lambda$ theorem)
\item[5.5] Theorem (uniqueness of measures)
\item[5.9] Theorem (monotone convergence)
\item[5.11] Theorem (Fatou's lemma)
\item[5.13] Theorem (dominated convergence)
\item[5.18] Proposition (inverse transform)
\item[5.20] Corollary (density of $g(X)$)
\item[5.24] Proposition (joint densities)
\item[5.25] Proposition (sums of independent Gaussians)
\item[5.26] Proposition (polar coordinates)
\item[5.27] Theorem (Box--Muller, 1958)
\end{results}
\subsection*{6.\ Convergence and the Central Limit Theorem}
\begin{results}
\item[6.4] Theorem (strong law, fourth moment)
\item[6.5] Theorem (Glivenko--Cantelli)
\item[6.6] Lemma (differentiation under the integral)
\item[6.8] Theorem (uniqueness and inversion)
\item[6.9] Corollary (averaging Cauchy variables)
\item[6.10] Theorem (Lindeberg, 1922)
\item[6.11] Corollary (central limit theorem)
\item[6.12] Theorem (local de Moivre--Laplace)
\item[6.13] Theorem (Stirling's constant from the CLT)
\end{results}
\subsection*{7.\ Random Walks}
\begin{results}
\item[7.2] Lemma (first-return decomposition)
\item[7.3] Theorem (the main lemma)
\item[7.5] Proposition (the maximum)
\item[7.6] Theorem (discrete arcsine law)
\item[7.7] Theorem (arcsine law)
\item[7.9] Theorem (gambler's ruin)
\item[7.12] Theorem (recurrence criterion)
\item[7.13] Theorem (Pólya, 1921)
\item[7.14] Theorem (discrete Dirichlet problem)
\item[7.15] Theorem (Papadimitriou)
\end{results}
\subsection*{8.\ Markov Chains}
\begin{results}
\item[8.2] Proposition (Markov property)
\item[8.8] Theorem (existence and uniqueness)
\item[8.9] Lemma (coupling inequality)
\item[8.11] Theorem (convergence theorem)
\item[8.13] Theorem (ergodic theorem)
\item[8.14] Proposition (detailed balance)
\item[8.16] Proposition (spectral bound)
\item[8.20] Theorem (Broder)
\item[8.21] Proposition (lower bound)
\item[8.23] Theorem (Bayer--Diaconis)
\item[8.24] Proposition (composition)
\end{results}
\subsection*{9.\ Poisson Processes}
\begin{results}
\item[9.1] Proposition (memorylessness)
\item[9.2] Lemma (gaps and arrival times)
\item[9.4] Theorem (uniform points)
\item[9.5] Lemma (colouring)
\item[9.6] Theorem (independent Poisson increments)
\item[9.7] Theorem (three definitions)
\item[9.8] Theorem (thinning)
\item[9.9] Corollary (superposition)
\item[9.11] Theorem (Le Cam's inequality)
\item[9.12] Theorem (M/M/1)
\end{results}
\subsection*{10.\ Conditional Expectation and Martingales}
\begin{results}
\item[10.1] Theorem (Riesz--Fischer)
\item[10.2] Theorem (orthogonal projection)
\item[10.4] Theorem (conditional expectation)
\item[10.9] Proposition (no system beats a fair game)
\item[10.12] Theorem (optional stopping, Doob)
\item[10.14] Theorem (waiting for a word)
\item[10.15] Theorem (Doob's maximal inequality)
\item[10.16] Corollary (Kolmogorov's inequality)
\item[10.18] Theorem (convergence in $L^2$)
\item[10.19] Lemma (Doob's upcrossing inequality)
\item[10.20] Theorem (martingale convergence, Doob)
\item[10.21] Theorem (P\'olya's urn)
\item[10.22] Theorem (growth of a branching process)
\item[10.24] Theorem (Azuma--Hoeffding)
\item[10.25] Theorem (bounded differences)
\end{results}
\subsection*{11.\ Brownian Motion}
\begin{results}
\item[11.4] Theorem (Wiener; L\'evy's construction)
\item[11.7] Proposition (time inversion)
\item[11.8] Theorem (quadratic variation)
\item[11.10] Theorem (Paley, Wiener and Zygmund, 1933)
\item[11.12] Theorem (reflection principle)
\item[11.14] Proposition (discrete It\^o formula)
\item[11.15] Proposition (one period)
\item[11.16] Theorem (binomial pricing)
\item[11.17] Theorem (Black--Scholes as a limit)
\end{results}
\subsection*{12.\ The Probabilistic Method}
\begin{results}
\item[12.2] Theorem (Erd\H os, 1947)
\item[12.3] Proposition (large cuts)
\item[12.5] Lemma (second moment)
\item[12.6] Theorem (threshold for triangles)
\item[12.7] Theorem (Lov\'asz local lemma, symmetric form)
\item[12.9] Theorem (small components)
\end{results}
\section{Tome 4: The Art of Calculation}
\subsection*{1.\ Finite Calculus}
\begin{results}
\item[1.2] Theorem (fundamental theorem of finite calculus)
\item[1.5] Proposition (summation by parts)
\item[1.7] Proposition (Abel summation)
\item[1.8] Theorem (Dirichlet's test)
\item[1.15] Theorem (Faulhaber's formula)
\item[1.16] Theorem (Newton's interpolation formula)
\end{results}
\subsection*{2.\ Generating Functions as a Computing Tool}
\begin{results}
\item[2.3] Lemma (summation by packets)
\item[2.4] Lemma (uniform convergence)
\item[2.9] Theorem (exponential formula)
\item[2.11] Proposition (Euler)
\item[2.12] Theorem (Euler's pentagonal number theorem)
\item[2.13] Corollary (Euler's recurrence)
\end{results}
\subsection*{3.\ Classic Series and Infinite Products}
\begin{results}
\item[3.4] Theorem (Euler)
\item[3.5] Theorem (Euler)
\item[3.6] Corollary (Wallis, 1656; Vi\`ete, 1593)
\item[3.7] Theorem (Euler, 1735--1740)
\end{results}
\subsection*{4.\ Integration Techniques}
\begin{results}
\item[4.1] Lemma (Heine)
\item[4.3] Theorem (Fubini on a rectangle)
\item[4.4] Theorem (domination)
\item[4.5] Corollary (dominated convergence for continuous functions)
\item[4.8] Theorem (Frullani)
\item[4.11] Theorem (Dirichlet)
\end{results}
\subsection*{5.\ The Gamma and Beta Functions}
\begin{results}
\item[5.6] Theorem (Bohr--Mollerup)
\item[5.7] Corollary (Weierstrass's product)
\item[5.11] Theorem (Euler)
\item[5.13] Theorem (Euler's reflection formula)
\item[5.14] Theorem (Legendre's duplication formula)
\end{results}
\subsection*{6.\ Fourier Series}
\begin{results}
\item[6.2] Proposition (best approximation)
\item[6.3] Corollary (Bessel's inequality and Riemann--Lebesgue)
\item[6.4] Lemma (Dirichlet kernel)
\item[6.5] Theorem (Dirichlet, 1829)
\item[6.8] Lemma (Fej\'er kernel)
\item[6.9] Theorem (Fej\'er, 1900)
\item[6.11] Theorem (Parseval)
\item[6.12] Theorem (Gibbs phenomenon)
\item[6.13] Theorem (Weyl's criterion, 1916)
\item[6.14] Corollary (Bohl, Sierpi\'nski, Weyl, 1909--1910)
\item[6.15] Proposition (Fourier inversion on $\Z^d$)
\item[6.17] Theorem (isoperimetric inequality)
\end{results}
\subsection*{7.\ The Discrete Fourier Transform and the FFT}
\begin{results}
\item[7.3] Theorem (inversion and Parseval)
\item[7.4] Proposition (aliasing)
\item[7.6] Theorem (convolution theorem)
\item[7.7] Corollary (multiplying polynomials)
\item[7.8] Lemma (butterfly)
\item[7.9] Theorem (Cooley--Tukey)
\item[7.11] Theorem (number-theoretic transform)
\end{results}
\subsection*{8.\ The Fourier Transform and Poisson Summation}
\begin{results}
\item[8.4] Lemma (Fubini on the plane)
\item[8.5] Lemma (Gaussian approximate identity)
\item[8.6] Theorem (Fourier inversion)
\item[8.7] Theorem (Plancherel)
\item[8.9] Theorem (Poisson summation)
\item[8.10] Corollary (theta transformation)
\item[8.12] Corollary (Theorem~3.4 again)
\item[8.13] Lemma (a Fresnel integral)
\item[8.14] Theorem (Gauss, 1805)
\item[8.15] Theorem (Shannon--Whittaker--Kotelnikov)
\item[8.16] Theorem (Heisenberg's inequality)
\end{results}
\subsection*{9.\ Complex Analysis for Computing}
\begin{results}
\item[9.2] Proposition (Cauchy--Riemann equations)
\item[9.3] Lemma (Abel)
\item[9.7] Theorem (Goursat)
\item[9.9] Theorem (Cauchy's theorem)
\item[9.10] Theorem (Cauchy's integral formula)
\item[9.11] Theorem (holomorphic means analytic)
\item[9.12] Corollary (Cauchy's estimates)
\item[9.14] Theorem (Morera)
\item[9.15] Theorem (limits and integrals)
\item[9.16] Theorem (Liouville)
\item[9.17] Corollary (fundamental theorem of algebra)
\item[9.18] Theorem (zeros and the identity theorem)
\item[9.19] Theorem (maximum principle)
\item[9.20] Theorem (Riemann's removable singularity theorem)
\item[9.22] Theorem (Laurent)
\item[9.24] Proposition (computing residues)
\item[9.26] Theorem (residue theorem)
\item[9.34] Corollary (Basel, the sixth proof)
\item[9.35] Corollary (Euler's partial fractions, on $\C$)
\item[9.37] Theorem (argument principle)
\item[9.38] Theorem (Rouch\'e)
\end{results}
\subsection*{10.\ The Zeta Function}
\begin{results}
\item[10.2] Proposition (products, uniqueness)
\item[10.3] Theorem (Euler product)
\item[10.5] Corollary (Euler, 1737)
\item[10.8] Theorem (Riemann, 1859)
\end{results}
\subsection*{11.\ Asymptotics}
\begin{results}
\item[11.3] Proposition (Euler's series)
\item[11.5] Theorem (Euler--Maclaurin formula)
\item[11.6] Corollary (harmonic numbers)
\item[11.8] Theorem (Stirling's series)
\item[11.9] Theorem (Laplace's method)
\item[11.13] Theorem (transfer)
\end{results}
\subsection*{12.\ The Arithmetic--Geometric Mean}
\begin{results}
\item[12.2] Theorem (Gauss, 1799)
\item[12.3] Corollary (Landen's transformation)
\item[12.4] Corollary (the lemniscate)
\item[12.7] Theorem (the perimeter by the AGM)
\item[12.9] Theorem (Legendre's relation, 1811)
\item[12.10] Theorem (Brent--Salamin, 1976)
\end{results}
\section{Tome 5: Where It All Meets}
\subsection*{1.\ Mertens's Theorems}
\begin{results}
\item[1.1] Lemma (Abel summation)
\item[1.2] Proposition (Chebyshev's bounds)
\item[1.3] Theorem (Mertens, 1874)
\item[1.4] Corollary (Chebyshev, 1852)
\item[1.5] Theorem (Mertens, 1874)
\item[1.6] Theorem (Mertens, 1874)
\end{results}
\subsection*{2.\ Arithmetic Functions on Average}
\begin{results}
\item[2.2] Theorem (Wintner, 1943)
\item[2.8] Lemma (Dirichlet's hyperbola method)
\end{results}
\subsection*{3.\ Dirichlet Characters and \texorpdfstring{$L$}{L}-Functions}
\begin{results}
\item[3.1] Proposition (orthogonality)
\item[3.8] Theorem (Dirichlet, 1837)
\item[3.10] Theorem (Dirichlet's formula)
\end{results}
\subsection*{4.\ Dirichlet's Theorem}
\begin{results}
\item[4.1] Theorem (Dirichlet, 1837)
\end{results}
\subsection*{5.\ The Prime Number Theorem}
\begin{results}
\item[5.2] Theorem (Hadamard, de la Vallée Poussin, 1896)
\item[5.3] Theorem (Newman, 1980)
\item[5.5] Theorem (the prime number theorem)
\item[5.7] Theorem (von Mangoldt, 1897)
\item[5.8] Theorem (the prime number theorem for arithmetic progressions)
\end{results}
\subsection*{6.\ The Zeros of \texorpdfstring{$\zeta$}{zeta}}
\begin{results}
\item[6.1] Theorem (Stirling's formula)
\item[6.3] Theorem (Jensen's formula)
\item[6.5] Lemma (Borel--Carathéodory)
\item[6.7] Theorem (Hadamard, 1893)
\item[6.9] Theorem (Riemann 1859, von Mangoldt 1905)
\item[6.12] Theorem (de la Vallée Poussin, 1899)
\end{results}
\subsection*{7.\ The Explicit Formula and the Riemann Hypothesis}
\begin{results}
\item[7.1] Lemma (Perron)
\item[7.2] Proposition (truncated Perron formula)
\item[7.3] Lemma (good heights)
\item[7.4] Lemma (the left half-plane)
\item[7.5] Theorem (the explicit formula, von Mangoldt 1895)
\item[7.6] Corollary (the exact formula)
\item[7.8] Theorem (de la Vallée Poussin, 1899)
\item[7.10] Theorem (von Koch, 1901)
\end{results}
\subsection*{8.\ Sieves}
\begin{results}
\item[8.1] Proposition (Legendre, 1808)
\item[8.2] Theorem (Selberg, 1947)
\item[8.4] Theorem (Brun--Titchmarsh)
\item[8.5] Theorem (Brun, 1919)
\item[8.6] Corollary (Brun, 1919)
\item[8.7] Lemma (Gallagher)
\item[8.8] Theorem (the large sieve inequality)
\item[8.10] Theorem (the arithmetic large sieve, Montgomery 1968)
\end{results}
\subsection*{9.\ The Erd\H{o}s--Kac Theorem}
\begin{results}
\item[9.1] Theorem (the Turán--Kubilius inequality)
\item[9.2] Theorem (Hardy and Ramanujan 1917, Turán 1934)
\item[9.7] Theorem (Erd\H{o}s and Kac, 1940)
\end{results}
\subsection*{10.\ Random Models of the Primes}
\begin{results}
\item[10.2] Theorem (Cramér)
\item[10.4] Theorem (Cramér, 1936)
\item[10.6] Proposition (a model that knows the small primes)
\item[10.7] Theorem (Gallagher, 1976)
\end{results}
\subsection*{11.\ Exponential Sums}
\begin{results}
\item[11.2] Lemma (van der Corput's inequality)
\item[11.3] Theorem (van der Corput's difference theorem)
\item[11.4] Theorem (Weyl, 1916)
\item[11.7] Theorem (Weyl's inequality for squares)
\item[11.8] Lemma (Kusmin--Landau)
\item[11.9] Theorem (van der Corput's second derivative test)
\item[11.10] Theorem (Fejér, van der Corput)
\item[11.11] Theorem (Pólya, Vinogradov, 1918)
\item[11.12] Theorem (Vinogradov, 1918)
\end{results}
\subsection*{12.\ A Taste of Modular Forms}
\begin{results}
\item[12.6] Theorem (Fourier expansion of $G_k$)
\item[12.7] Theorem (Dedekind)
\item[12.13] Theorem (Hecke)
\item[12.14] Theorem (Mordell, 1917)
\end{results}
